\documentclass[11pt]{article}
\usepackage[T1]{fontenc}
\usepackage[utf8]{inputenc}
\usepackage{lmodern,amsmath,amssymb,amsthm,booktabs,array}
\usepackage[margin=0.85in]{geometry}
\usepackage[colorlinks=true,urlcolor=blue,citecolor=blue,linkcolor=blue]{hyperref}
\setlength{\emergencystretch}{1em}
\newtheorem{theorem}{Theorem}
\newtheorem{lemma}[theorem]{Lemma}
\newtheorem{corollary}[theorem]{Corollary}
\newcommand{\Good}{\operatorname{Good}}
\newcommand{\Conj}{\operatorname{Conj}}
\newcommand{\Cl}{\operatorname{Cl}}
\newcommand{\Inv}{\operatorname{Inv}}
\newcommand{\Aut}{\operatorname{Aut}}
\title{Negative answers to Ta's Problems 11.6--11.7\\
\large arXiv:2505.08090v5, Section 11, page 27;\\
and a proof of Conjecture 5.13\\
arXiv:2508.16772v1, Section 5.2.2, page 17:\\
central-twist counting extending Kamada--Oshiro (2009),\\
Ta (2025), and Nakajima--Yamaguchi (2026)}
\author{Haobo Ma\\{\small ChronoAI Pte Ltd, Singapore}
\and Wenlin Zhang\\{\small National University of Singapore}\\
{\small ChronoAI Pte Ltd, Singapore}}
\date{November 9, 2026}
\begin{document}
\maketitle
\begin{abstract}
We answer negatively Ta's Problems 11.6--11.7 in arXiv:2505.08090v5,
Section 11, page 27, and prove his Conjecture 5.13 in
arXiv:2508.16772v1, Section 5.2.2, page 17, extending the counting
methods of Kamada--Oshiro (2009), Ta (2025), and
Nakajima--Yamaguchi (2026).
The two semidirect products $C_{16}\rtimes C_4$ in which a generator of
$C_4$ acts by multiplication by $5$ or $7$ both have order $64$, center
order $4$, and $22$ conjugacy classes, but their conjugation quandles have
$8\,847\,360$ and $2\,160\,000$ good involutions, respectively.
Ta's central-label characterization gives a weighted-involution formula
for every finite conjugation subquandle, including an exact criterion
for equality in the central-label bound. The quaternion example attains
that bound with two orbits and a center of order two. We also distinguish
inversion-closure from classwise inverse stability, correcting Remark 2.14.
For the linear quandle $\Lambda_{4n,2n-1}$, the number of good
involutions is ten for odd $n$ and four for even $n$; the odd case
proves Conjecture 5.13, and the even case reproves Ta's Remark 5.14.
\end{abstract}
\begin{center}\small
\copyright\ 2026 Haobo Ma and Wenlin Zhang.
Licensed under Creative Commons Attribution 4.0 International (CC BY 4.0).
\end{center}

\section{The question and the two groups}
For a group $G$, put $s_x(y)=xyx^{-1}$. Its conjugation quandle is
$\Conj G=(G,s)$. A \emph{good involution} is a map $\rho:G\to G$ such that
\begin{equation}\label{eq:good}
 \rho^2=\mathrm{id},\qquad \rho s_x=s_x\rho,\qquad
 s_{\rho(x)}=s_x^{-1}\quad(x\in G).
\end{equation}
An involution here may have fixed points. We count maps, not isomorphism
classes of symmetric quandles. Ta's Problem 11.6 \cite[p.~27]{Ta}
poses the determination of $|\Good(\Conj G)|$ by
$(|G|,|Z(G)|,k(G))$, where $k(G)=|\Cl(G)|$, for finite nonabelian groups.
Its conditional second part requests an exact formula in those three numbers.
The source asks whether the count is ``completely determined'' by the
triple; Theorem~\ref{thm:main} gives a negative answer.
The questions occur already in version 3 (July 15, 2025), whereas
versions 1--2 concern Legendrian racks and do not contain these questions.

For $r\in\{5,7\}$, define $G_r$ on
$(\mathbb Z/16)\times(\mathbb Z/4)$ by
\begin{equation}\label{eq:product}
 (i,j)(k,l)=(i+r^j k,j+l).
\end{equation}
The first coordinate is reduced modulo $16$ and the second modulo $4$.
Since $r^4\equiv1\pmod {16}$ and $r$ is a unit, this is a semidirect product.
Write $a=(1,0)$ and $b=(0,1)$, so $bab^{-1}=a^r$.
In particular, both groups are nonabelian.

\begin{samepage}
\begin{theorem}\label{thm:main}
For $r=5,7$ one has
\[
 (|G_r|,|Z(G_r)|,k(G_r))=(64,4,22),
\]
whereas
\[
 |\Good(\Conj G_5)|=10\cdot4^6\cdot6^3=8\,847\,360,
 \qquad
 |\Good(\Conj G_7)|=10^4\cdot6^3=2\,160\,000.
\]
Thus neither determination nor an exact formula depending only on the
three numbers in Problem 11.6 is possible.
\end{theorem}
\end{samepage}

The central-label description used below is Ta's Theorem 6.7, with the
label $z$ equal to the inverse of Ta's $\zeta$ \cite[pp.~13--14]{Ta}.
We include the proof for finite conjugation subquandles and the counting
reduction. The central factor of the formula recovers the specialization
of Nakajima--Yamaguchi's decomposition theorem
\cite[Theorems 1.1 and 3.3]{NY}; the quaternion subquandle below is
the dihedral quandle $R_4$, whose four good involutions were classified
by Kamada and Oshiro \cite[Theorem 3.2(3)]{KO}.
Ta's subsequent work already counts weighted involutions of an orbit set
and factors the choices over parity components for the Alexander family
$\Lambda_{4n,2n+1}$ \cite[Lemma 5.10 and Theorem 5.11, pp.~15--16]{TaTwisted}.
The formula below extends that counting mechanism to arbitrary weights
from central twisted inversion on finite conjugation subquandles.
Corollary~\ref{cor:linear} also proves the cross-referenced Conjecture 5.13
for $\Lambda_{4n,2n-1}$; its even companion is already Ta's Remark 5.14.

\section{An exact formula from central twists}
Let $X$ be a nonempty conjugation-closed subset of a finite group, let
$H=\langle X\rangle$, and let $\mathcal C=X/H$ be its set of
$H$-conjugacy classes. Thus conjugation by $H$ preserves $X$.
For $X=G$ this gives $H=G$ and $\mathcal C=\Cl(G)$.
\begin{lemma}\label{lem:labels}
Good involutions of $\Conj X$ correspond bijectively to functions
$z:X\to Z(H)$ that are constant on $H$-conjugacy classes and satisfy
\begin{equation}\label{eq:labels}
 z(x)x^{-1}\in X,\qquad z\bigl(z(x)x^{-1}\bigr)=z(x)\quad(x\in X).
\end{equation}
The corresponding map is $\rho(x)=z(x)x^{-1}$.
\end{lemma}
\begin{proof}
If $\rho$ satisfies \eqref{eq:good} on $X$, then
$s_{\rho(x)x}=s_{\rho(x)}s_x=\mathrm{id}$ on $X$. Hence
$\rho(x)x$ centralizes $X$ and therefore $H$, so
$z(x):=\rho(x)x\in Z(H)$. Commutation with all $s_y$, $y\in X$, gives
\[
 z(yxy^{-1})=y\rho(x)xy^{-1}=z(x).
\]
This invariance extends to $H$, since $X$ generates $H$ and invariance
under a bijection implies invariance under its inverse.
Since $\rho(x)=z(x)x^{-1}$ and $z(x)$ is central,
$\rho^2(x)=z(\rho(x))xz(x)^{-1}$. Thus $\rho^2=\mathrm{id}$ is
equivalent to \eqref{eq:labels}. Conversely, a class-constant central
function satisfying \eqref{eq:labels} defines an involution by that same
identity. Class constancy gives $\rho s_y=s_y\rho$, and centrality gives
$s_{\rho(x)}=s_{x^{-1}}=s_x^{-1}$. Finally, $z(x)=\rho(x)x$ proves uniqueness.
\end{proof}

For $C,D\in\mathcal C$, define
\begin{equation}\label{eq:weights}
 T_z(C)=zC^{-1}\quad(z\in Z(H)),\qquad
 w_{CD}=|\{z\in Z(H):T_z(C)=D\}|.
\end{equation}
These are well-defined $H$-conjugacy classes, and $T_z^2=\mathrm{id}$
on the classes of $H$. A class $T_z(C)$ outside $X$ contributes to no
$w_{CD}$. Thus every row sum is at most $|Z(H)|$.
Consequently $w_{CD}=w_{DC}$. Let $\Inv(S)$ be the permutations of a finite
set $S$ whose square is the identity, and let $\mathcal P(\sigma)$ denote
the unordered two-cycles of such a permutation.

\begin{corollary}\label{cor:formula}
For every finite conjugation subquandle as above,
\begin{equation}\label{eq:formula}
 |\Good(\Conj X)|=
 \sum_{\sigma\in\Inv(\mathcal C)}
 \prod_{\sigma(C)=C}w_{CC}
 \prod_{\{C,D\}\in\mathcal P(\sigma)}w_{CD}.
\end{equation}
The count factors over the connected components of the graph with vertex
set $\mathcal C$ and an edge between distinct $C,D$ when $w_{CD}>0$.
\end{corollary}
\begin{proof}
By Lemma~\ref{lem:labels}, a good involution induces an involution $\sigma$
of the class set. At a fixed class $C$, its label has $w_{CC}$ choices.
For a pair $C,D$ exchanged by $\sigma$, equation \eqref{eq:labels}
requires the same central label on both classes. Exactly $w_{CD}$ labels
work, since $T_z(C)=D$ also implies $T_z(D)=C$.
These choices are independent across cycles, and every choice satisfies
\eqref{eq:labels}. This proves both completeness and the formula.
A nonzero summand cannot pair classes in different graph components,
so the choices also factor over those components.
\end{proof}

The summation over involutions and factorization in
Corollary~\ref{cor:formula} have the precedent just cited in
Ta \cite[Lemma 5.10 and Theorem 5.11]{TaTwisted}.
There the orbit choices in a linear Alexander family have a uniform
two-choice weight and split by parity. Here the weights $w_{CD}$ depend
on the conjugacy classes and may vanish or vary; the graph records exactly
which components can be counted independently.

\begin{corollary}\label{cor:equality}
Put $q=|Z(H)|$ and $k=|\mathcal C|$. If $q\ge2$, then
\[
 |\Good(\Conj X)|=q^k
 \quad\Longleftrightarrow\quad
 zC^{-1}=C\quad(C\in\mathcal C,\ z\in Z(H)).
\]
For $q=1$, equality holds exactly when $X=X^{-1}$.
\end{corollary}
\begin{proof}
There are $q^k$ class-constant label functions. If every central
twisted inversion preserves each class, all these functions satisfy
Lemma~\ref{lem:labels}. Conversely, suppose all functions give good maps.
A label sending a class outside $X$ is impossible. If $zC^{-1}=D\ne C$,
choose label $z$ on $C$ and a different label on $D$; this violates
\eqref{eq:labels}. Thus every twist preserves every class.
If $q=1$, the sole label is $1$, giving precisely inversion on $X$.
\end{proof}

\begin{corollary}\label{cor:center}
For a finite group $G$, put $q=|Z(G)|$, $k=k(G)$, and
\[
 I(q)=\sum_{j=0}^{\lfloor q/2\rfloor}
       \frac{q!}{2^j j!(q-2j)!}.
\]
Then
\[
 |\Good(\Conj G)|\le I(q)q^{k-q}.
\]
In particular, $|\Good(\Conj G)|<q^k$ whenever $q\ge2$.
The sharper bound is attained by $G=Q_8$.
\end{corollary}
\begin{proof}
The $q$ central singleton classes form a separate block with matrix
$J_q$: for $c,d\in Z(G)$, exactly one label, $dc$, sends $c$ to $d$.
Central twists cannot send a noncentral element into the center.
Consequently this block contributes $I(q)$, the number of involutions
on $q$ points; choosing $j$ disjoint pairs gives the displayed sum.
The other $k-q$ classes have at most $q$ labels each. This proves the
bound, and $I(q)<q^q$ for $q\ge2$ because not every function on $q$
points is an involutive permutation.
In $Q_8$, the classes are $\{1\}$, $\{-1\}$,
$\{\pm i\}$, $\{\pm j\}$, and $\{\pm k\}$.
The last three classes are preserved by both central twists of
inversion, giving $I(2)2^3=16$.
\end{proof}

The central-block factor $I(q)$ is the group case of
Nakajima--Yamaguchi \cite[Theorems 1.1 and 3.3]{NY}: the maximal
trivial component of $\Conj G$ is $Z(G)$, and its good involutions
are all involutions of that set. The bound above and the classwise
criterion specify the remaining central-label restrictions.

\section{The class calculations and the two counts}
In the coordinates of \eqref{eq:product}, direct multiplication gives
\begin{align}
 (i,j)^{-1}&=(-r^{-j}i,-j),\label{eq:inverse}\\
 (u,v)(i,j)(u,v)^{-1}&=((1-r^j)u+r^v i,j).\label{eq:conjugacy}
\end{align}
Commuting with $a$ and $b$ shows that $(i,j)$ is central precisely when
$r^j\equiv1$ and $(r-1)i\equiv0$ modulo $16$. Hence
\[
 Z(G_5)=\{(0,0),(4,0),(8,0),(12,0)\},\qquad
 Z(G_7)=\{(0,0),(0,2),(8,0),(8,2)\}.
\]

Let $J_m$ denote the $m\times m$ all-one matrix and put
$P=\left(\begin{smallmatrix}0&4\\4&0\end{smallmatrix}\right)$.
Tables~\ref{tab:five} and \ref{tab:seven} give every block of the
weight matrix \eqref{eq:weights}. Each listed pair represents a distinct
conjugacy class; the sizes are in the same order. Thus the tables also
specify all conjugacy classes, not just selected examples.

\begin{table}[htbp]\centering\small
\begin{tabular}{@{}lll r@{}}\toprule
Class representatives in $G_5$ & Class sizes & Weight matrix & Factor\\\midrule
$(0,0),(4,0),(8,0),(12,0)$ & $1,1,1,1$ & $J_4$ & $10$\\
$(0,1),(0,3)$ & $4,4$ & $P$ & $4$\\
$(0,2),(4,2)$ & $2,2$ & $2J_2$ & $6$\\
$(1,0),(3,0)$ & $4,4$ & $P$ & $4$\\
$(1,1),(3,3)$ & $4,4$ & $P$ & $4$\\
$(1,2),(3,2)$ & $4,4$ & $P$ & $4$\\
$(1,3),(3,1)$ & $4,4$ & $P$ & $4$\\
$(2,0),(6,0)$ & $2,2$ & $2J_2$ & $6$\\
$(2,1),(2,3)$ & $4,4$ & $P$ & $4$\\
$(2,2),(6,2)$ & $2,2$ & $2J_2$ & $6$\\\bottomrule
\end{tabular}
\caption{The ten central-twist blocks for $G_5$.}\label{tab:five}
\end{table}

\begin{table}[htbp]\centering\small
\begin{tabular}{@{}lll r@{}}\toprule
Class representatives in $G_7$ & Class sizes & Weight matrix & Factor\\\midrule
$(0,0),(0,2),(8,0),(8,2)$ & $1,1,1,1$ & $J_4$ & $10$\\
$(0,1),(0,3)$ & $8,8$ & $2J_2$ & $6$\\
$(1,0),(1,2),(9,0),(9,2)$ & $2,2,2,2$ & $J_4$ & $10$\\
$(1,1),(1,3)$ & $8,8$ & $2J_2$ & $6$\\
$(2,0),(2,2),(6,0),(6,2)$ & $2,2,2,2$ & $J_4$ & $10$\\
$(3,0),(3,2),(11,0),(11,2)$ & $2,2,2,2$ & $J_4$ & $10$\\
$(4,0),(4,2)$ & $2,2$ & $2J_2$ & $6$\\\bottomrule
\end{tabular}
\caption{The seven central-twist blocks for $G_7$.}\label{tab:seven}
\end{table}

Here is an explicit way to recover each entry. For a representative
$(i,j)$, equation \eqref{eq:conjugacy} gives exactly the class
\begin{equation}\label{eq:classset}
 \{((1-r^j)u+r^v i,j):0\leq u<16,\ 0\leq v<4\}.
\end{equation}
For $z=(c,d)$ in the displayed center, equations \eqref{eq:product}
and \eqref{eq:inverse} give
\begin{equation}\label{eq:twist}
 z(i,j)^{-1}=(c-r^{d-j}i,d-j).
\end{equation}
Locating this element in \eqref{eq:classset} for each of the four $z$'s
gives the corresponding row of the weight matrix.
For instance, in $G_5$ the classes of $(0,1)$ and $(0,3)$ have first
coordinate divisible by $4$; all four central twists exchange them.
For $(0,2)$ and $(4,2)$ the classes are respectively
$\{(0,2),(8,2)\}$ and $\{(4,2),(12,2)\}$, and two labels lead to
each target. In $G_7$ the odd-second-coordinate classes are determined
by parity of the first coordinate. The choices $d=0,2$ interchange or
preserve the second coordinate, while $c=0,8$ give the same parity;
this yields the two $2J_2$ blocks of class size $8$.
The other entries follow by the same displayed formulas, with all their
representatives and sizes supplied in the tables. Their class sets are
pairwise disjoint and their sizes sum to $64$ for each group, establishing
that the lists are exhaustive. There are $22$ representatives in each list.

For a $J_4$ block, formula \eqref{eq:formula} counts the involutions on
four objects: $1+\binom42+3=10$. A $P$ block has only the transposition,
with four labels. A $2J_2$ block has $2\cdot2$ choices for two fixed
classes and two choices for the transposition, giving $6$.
Multiplying the factors in the two tables proves Theorem~\ref{thm:main}.

\section{The bound question and inverse classes}
Problem 11.7 \cite[Section 11, p.~27]{Ta} asks whether the upper
bound in Corollary 6.12 is ``always strict'' when both the center
and the orbit set have size at least two. The distinction between
a whole group and a proper subquandle is essential.

\begin{corollary}\label{cor:quaternion}
Let $X=\{\pm i,\pm j\}\subset Q_8$ with its conjugation quandle.
Then $H=\langle X\rangle=Q_8$, $|Z(H)|=2$, $|X/H|=2$, and
\[
 |\Good(\Conj X)|=4
 =\min\bigl(|\Aut(\Conj X)|,\,2^2\bigr),
 \qquad |\Aut(\Conj X)|=8.
\]
Thus Problem 11.7 has a negative answer. Moreover,
$\operatorname{Anti}(\Conj X)=\Aut(\Conj X)\cong C_2\wr S_2$,
and exactly three symmetric-quandle isomorphism classes arise from $X$.
\end{corollary}
\begin{proof}
The two $H$-orbits are $A=\{\pm i\}$ and $B=\{\pm j\}$.
The permutations $s_i=s_{-i}$ fix $A$ and exchange the two elements
of $B$; $s_j=s_{-j}$ do the reverse. Inversion and multiplication by
$-1$ each preserve both orbits. Corollary~\ref{cor:equality} therefore
gives four good involutions: independently exchange or fix each pair.
Any automorphism permutes the two inner orbits, and conversely every
permutation preserving the unordered partition $\{A,B\}$ respects
the displayed operations. These are the eight permutations of
$C_2\wr S_2$. Since every $s_x$ is an involution, the quandle equals
its dual, so its antiautomorphisms are the same eight permutations.
Conjugating the four good maps by these automorphisms gives three
orbits: the identity, exchange of both pairs, and exchange of one pair.
These are exactly the symmetric-quandle isomorphism classes.
\end{proof}

Under the assignment $0\mapsto i$, $1\mapsto j$, $2\mapsto-i$,
$3\mapsto-j$, the operation is $s_x(y)=2x-y\pmod4$.
Thus $\Conj X\cong R_4$, and its four good involutions reprove the
$n=4$ case of Kamada--Oshiro \cite[Theorem 3.2(3)]{KO}, released
in 2009 and published in 2010. Their theorem supplies the prior
classification; the realization inside $Q_8$ applies it to the two
hypotheses and the full minimum bound in Ta's later Problem 11.7.

\begin{corollary}\label{cor:inverse}
For a finite conjugation subquandle $X$, inversion-closure means that
inversion permutes the $H$-classes contained in $X$, not necessarily
that it fixes each class. In particular, the classwise equivalence
in Ta's Remark 2.14 (Section 2.1.1, p.~6) requires this correction.
\end{corollary}
\begin{proof}
Inversion sends the class of $x$ to the class of $x^{-1}$ and has
square the identity. Hence $X=X^{-1}$ is equivalent to closure of
the set of classes under that permutation. For an explicit distinction,
take $X=H=G_5$. It is inversion-closed, but
\[
 a^{G_5}=\{a,a^5,a^9,a^{13}\},\qquad
 (a^{-1})^{G_5}=\{a^3,a^7,a^{11},a^{15}\}
\]
are different classes by \eqref{eq:conjugacy}.
Thus $a$ is not ``conjugate to its inverse'', the condition asserted
in that remark. In the weighted formula the inverse classes are
allowed to be paired, rather than required to coincide.
\end{proof}

\section{A cross-referenced linear family}
For $n\ge1$, put $m=4n$ and $k=2n-1$, and let
$\Lambda_{4n,2n-1}$ be the linear quandle on $\mathbb Z/m$ with
$s_x(y)=ky+(1-k)x$. Since $k^2\equiv1\pmod m$, every $s_x$
is an involution. Ta's Conjecture 5.13
\cite[Section 5.2.2, p.~17]{TaTwisted} asserts that
$|\Good\Lambda_{4n,2n-1}|=10$ for odd $n\ge3$.

\begin{corollary}\label{cor:linear}
For every positive integer $n$,
\[
 |\Good\Lambda_{4n,2n-1}|=
 \begin{cases}10,&n\text{ odd},\\4,&n\text{ even}.\end{cases}
\]
For odd $n$, Chinese remainder coordinates identify all good maps as
$(u,v)\mapsto(\sigma(u),v)$ on $(\mathbb Z/4)\times(\mathbb Z/n)$,
where $\sigma$ is an arbitrary involution of four points.
The odd case proves Conjecture 5.13; $n=1$ is an additional case not
posed there. The even case reproves Ta's Remark 5.14.
\end{corollary}
\begin{proof}
Suppose first that $n$ is odd. Modulo $4$ the coefficient $k$ is $1$,
and modulo $n$ it is $-1$. Thus the Chinese remainder isomorphism gives
\[
 s_{(a,b)}(u,v)=(u,2b-v).
\]
Write $\rho(u,v)=(f(u,v),g(u,v))$. The identity
$s_{\rho(u,v)}=s_{(u,v)}^{-1}=s_{(u,v)}$ forces $2g(u,v)=2v$
modulo $n$, hence $g(u,v)=v$. Commutation with all $s_{(a,b)}$
then gives $f(u,2b-v)=f(u,v)$. For fixed $v$, the values $2b-v$
exhaust $\mathbb Z/n$, so $f(u,v)=\sigma(u)$ is independent of $v$.
The remaining identity $\rho^2=\mathrm{id}$ says exactly that
$\sigma^2=\mathrm{id}$. Conversely, each such $\sigma$ satisfies
all three good-involution identities. There are $1+6+3=10$ involutions
on four points. The argument includes the trivial factor $n=1$.

For even $n$, the label description of Ta
\cite[Theorem 4.2]{TaTwisted} can be seen directly here.
For a good map put $\psi(x)=\rho(x)+kx$. Equality $s_{\rho(x)}=s_x$
implies $(1-k)\psi(x)=0$. Commutation gives
$\psi(s_a(x))=k\psi(x)=\psi(x)$, so $\psi$ is constant on inner
orbits. Finally $\rho^2(x)=x+\psi(\rho(x))-\psi(x)$.
Conversely, these label conditions imply all the good identities for
$\rho(x)=\psi(x)-kx$.
Since $\gcd(1-k,4n)=2$ for even $n$, the fixed-label set is
$\ker(1-k)=\{0,2n\}$. The products $s_a s_b$ are translations by
$(1-k)(a-b)$ and generate all even translations, while every $s_a$
preserves parity. Consequently the inner orbits are precisely the two
parity classes. Each assignment of one of $0,2n$ to each class defines
$\rho$ preserving parity, so $\psi(\rho(x))=\psi(x)$ automatically.
All four assignments work, and the label description is exhaustive.
\end{proof}

\section{Coverage and verification}
The weights in Corollary~\ref{cor:formula}, rather than the three
cardinalities in Problem 11.6, determine the count. The source coverage
is recorded below; a partial instance does not classify all racks.

\begin{center}\small
\begin{tabular}{@{}>{\raggedright\arraybackslash}p{0.23\linewidth}>{\raggedright\arraybackslash}p{0.70\linewidth}@{}}\toprule
Source item & Coverage\\\midrule
Problem 11.1, p.~27 & Partial answer: eight antiautomorphisms for $X$
in Corollary~\ref{cor:quaternion}; general families not covered.\\
Problem 11.2, p.~27 & Partial answer: $\Aut\cup\operatorname{Anti}
=C_2\wr S_2$ for this $X$; general families not covered.\\
Problem 11.3, p.~27 & Partial answer: the four good maps of $X$
(reproved; prior classification Kamada--Oshiro, 2009, $R_4$).
Non-conjugation families not covered.\\
Problem 11.4, p.~27 & Not covered: $n$-fold conjugation operations.\\
Problem 11.5, p.~27 & Partial answer: three symmetric-quandle
isomorphism classes for $X$; other subquandles not covered.\\
Problem 11.6, p.~27 & Answered: no, for determination and for a formula
in those three numbers (Theorem~\ref{thm:main}).\\
Problem 11.7, p.~27 & Answered: no (Corollary~\ref{cor:quaternion});
the full minimum bound equals $4$.\\
Conjecture 11.8, p.~27 & Prior proof cited: Nakajima--Yamaguchi
(2026), Example 6.8, pp.~24--26, give a connected noninvolutory
quandle with more than one good involution.\\
Conjecture 11.9, p.~28 & Prior proof cited: the same conjugation
example and following paragraph of \cite{NY}, p.~26.\\
Remark 2.14, p.~6 & Refuted-and-corrected:
Corollary~\ref{cor:inverse} replaces classwise inverse stability
by closure of the class set under inversion.\\
Theorem 6.7; Corollary 6.12, pp.~13--14 & Reproved specialization
of Ta (2025); weighted refinement and exact equality criterion in
Corollaries~\ref{cor:formula}--\ref{cor:equality}.\\
Remark 6.13, p.~14 & Extended: the proper-subquandle equality
in Corollary~\ref{cor:quaternion} contrasts with the strict whole-group
bound in Corollary~\ref{cor:center}.\\\bottomrule
\end{tabular}
\end{center}

\begin{center}\small
\begin{tabular}{@{}>{\raggedright\arraybackslash}p{0.23\linewidth}>{\raggedright\arraybackslash}p{0.70\linewidth}@{}}\toprule
Ta, arXiv:2508.16772v1 & Coverage\\\midrule
Lemma 5.10; Theorem 5.11, pp.~15--16 & Prior method cited:
weighted orbit involutions and parity factorization;
Corollary~\ref{cor:formula} extends the mechanism to central-twist weights.\\
Conjecture 5.13, Section 5.2.2, p.~17 & Proved:
$|\Good\Lambda_{4n,2n-1}|=10$ for every odd $n\ge3$
(Corollary~\ref{cor:linear}); $n=1$ is an extension not posed there.\\
Remark 5.14, p.~17 & Reproved; prior proof Ta (2025):
the even-$n$ count is four (Corollary~\ref{cor:linear}).\\\bottomrule
\end{tabular}
\end{center}

The written arguments establish Lemma~\ref{lem:labels},
Corollaries~\ref{cor:formula}--\ref{cor:center}, the class tables,
the two counts, and Corollaries~\ref{cor:quaternion}--\ref{cor:linear}.
The accompanying \texttt{verify.py} independently constructs both groups,
enumerates their centers and classes, and counts all compatible central
labels block by block. Its largest label enumeration has $4^4$ cases.
It also checks the good-involution identities on every accepted block
assignment. The output is \texttt{verification.json}.
The accompanying \texttt{verify\_siblings.py} checks all $4^4$ maps
of the quaternion subquandle directly, its eight automorphisms and
eight antiautomorphisms, the three symmetric isomorphism classes,
and the sharper whole-$Q_8$ bound. It also checks the weighted formula
and equality criterion on every nonempty conjugation-closed subset
of $Q_8$, with its own generated subgroup and center. For every
$n=1,\ldots,31$, it enumerates every allowed label assignment for
$\Lambda_{4n,2n-1}$, checks all defining good identities for each accepted
map, and confirms the ten/four parity count in Corollary~\ref{cor:linear}.
Its output is \texttt{sibling-verification.json}. These finite checks support the
examples; the general formulas and criteria have the written proofs above.

The pinned Lean file
\href{https://github.com/the-omega-institute/trureturing/blob/d9880e9f2c329bb3db982f33791419d5d90126b9/D5/S3/FiniteGroups/Quandles/TaThreeInvariantRefutation.lean}{\texttt{TaThreeInvariantRefutation.lean}}
contains \texttt{good\_counts}, proving both complete numerical counts,
and \texttt{result}, proving the negation of the finite-nonabelian
three-invariant determination assertion. Its \texttt{Gr false} and
\texttt{Gr true} are $G_5$ and $G_7$, and \texttt{goodCount} counts
all whole-group good maps. These statements concern the same maps as
\eqref{eq:good}, not their isomorphism classes. The companion
\href{https://github.com/the-omega-institute/trureturing/blob/d9880e9f2c329bb3db982f33791419d5d90126b9/D5/S3/FiniteGroups/Quandles/TaGoodInvolutionBoundRefutation.lean}{\texttt{TaGoodInvolutionBoundRefutation.lean}}
has \texttt{result}, negating the strictness assertion with precisely
the quaternion witness: four good maps, eight automorphisms, center
order two, and two orbits. The three symmetric isomorphism classes,
the general subquandle criterion, the correction to Remark 2.14,
and Corollary~\ref{cor:linear}
have the written and finite scopes just specified.

The shared
\href{https://github.com/the-omega-institute/trureturing/blob/d9880e9f2c329bb3db982f33791419d5d90126b9/D5/S3/FiniteGroups/Quandles/ConjugationSubquandle.lean}{\texttt{ConjugationSubquandle.lean}}
defines good maps, automorphisms, and \texttt{orbitCount}.
The whole-group label bijection is \texttt{goodLabelEquiv} in
\href{https://github.com/the-omega-institute/trureturing/blob/d9880e9f2c329bb3db982f33791419d5d90126b9/D5/S3/FiniteGroups/Quandles/CentralLabels.lean}{\texttt{CentralLabels.lean}}.
Its \texttt{wholeGoodEquiv} identifies those maps with the shared definition.
The complete block factorization appears as \texttt{labelBlockEquiv}
and \texttt{whole\_good\_card} in
\href{https://github.com/the-omega-institute/trureturing/blob/d9880e9f2c329bb3db982f33791419d5d90126b9/D5/S3/FiniteGroups/Quandles/BlockCounting.lean}{\texttt{BlockCounting.lean}}.
The declaration \texttt{classEquiv} identifies its partition with the
conjugacy-class quotient.

\begin{thebibliography}{9}\small
\bibitem{Ta}
L. Ta, \emph{Good involutions of conjugation subquandles},
arXiv:2505.08090v5, August 4, 2025.
\url{https://arxiv.org/abs/2505.08090v5}.

\bibitem{TaTwisted}
L. Ta, \emph{Good involutions of twisted conjugation subquandles and
Alexander quandles}, arXiv:2508.16772v1, August 22, 2025.
\url{https://arxiv.org/abs/2508.16772v1}.

\bibitem{NY}
Y. Nakajima and K. Yamaguchi, \emph{Decompositions of good involutions on
quandles}, arXiv:2608.19951v1, August 20, 2026.
\url{https://arxiv.org/abs/2608.19951v1}.

\bibitem{KO}
S. Kamada and K. Oshiro, \emph{Homology groups of symmetric quandles
and cocycle invariants of links and surface-links},
Trans. Amer. Math. Soc. \textbf{362} (2010), 5501--5527.
\url{https://doi.org/10.1090/S0002-9947-2010-05131-1}.
Earlier version: arXiv:0902.4277v1, February 25, 2009,
\url{https://arxiv.org/abs/0902.4277v1}.
\end{thebibliography}
\end{document}
